% ------------------------------------------------------------------------------------------------------------------
% Basic configuration and packages
% ------------------------------------------------------------------------------------------------------------------
% Package for discovering wrong and outdated usage of LaTeX.
% More information to be found in l2tabu English version.
\RequirePackage[l2tabu, orthodox]{nag}
% Class of LaTeX document: {size of paper, size of font}[document class]
\documentclass[a4paper,11pt]{article}



% ------------------------------------------------------
% Packages not tied to particular normal language
% ------------------------------------------------------
% This package should improved spaces in the text
\usepackage{microtype}
% Add few important symbols, like text Celcius degree
\usepackage{textcomp}



% ------------------------------------------------------
% Polonization of LaTeX document
% ------------------------------------------------------
% Basic polonization of the text
\usepackage[MeX]{polski}
% Switching on UTF-8 encoding
\usepackage[utf8]{inputenc}
% Adding font Latin Modern
\usepackage{lmodern}
% Package is need for fonts Latin Modern
\usepackage[T1]{fontenc}



% ------------------------------------------------------
% Setting margins
% ------------------------------------------------------
\usepackage[a4paper, total={14cm, 25cm}]{geometry}



% ------------------------------------------------------
% Setting vertical spaces in the text
% ------------------------------------------------------
% Setting space between lines
\renewcommand{\baselinestretch}{1.1}

% Setting space between lines in tables
\renewcommand{\arraystretch}{1.4}



% ------------------------------------------------------
% Packages for scientific papers
% ------------------------------------------------------
% Switching off \lll symbol, that I guess is representing letter "Ł"
% It collide with `amsmath' package's command with the same name
\let\lll\undefined
% Basic package from American Mathematical Society (AMS)
\usepackage[intlimits]{amsmath}
% Equations are numbered separately in every section
\numberwithin{equation}{section}

% Other very useful packages from AMS
\usepackage{amsfonts}
\usepackage{amssymb}
\usepackage{amscd}
\usepackage{amsthm}

% Package with better looking calligraphy fonts
\usepackage{calrsfs}

% Package with better looking greek letters
% Example of use: pi -> \uppi
\usepackage{upgreek}
% Improving look of lambda letter
\let\oldlambda\Lambda
\renewcommand{\lambda}{\uplambda}




% ------------------------------------------------------
% BibLaTeX
% ------------------------------------------------------
% Package biblatex, with biber as its backend, allow us to handle
% bibliography entries that use Unicode symbols outside ASCII
\usepackage[
language=polish,
backend=biber,
style=alphabetic,
url=false,
eprint=true,
]{biblatex}

\addbibresource{Wstęp-do-teorii-pola-Bibliography.bib}





% ------------------------------------------------------
% Defining new environments (?)
% ------------------------------------------------------
% Defining enviroment "Wniosek"
\newtheorem{corollary}{Wniosek}
\newtheorem{definition}{Definicja}
\newtheorem{theorem}{Twierdzenie}





% ------------------------------------------------------
% Local packages
% You need to put them in the same directory as .tex file
% ------------------------------------------------------
% Local configuration of this particular file
\usepackage{./Local-packages/local-settings}

% Package containing various command useful for working with a text
\usepackage{./Local-packages/general-commands}
% Package containing commands and other code useful for working with
% mathematical text
\usepackage{./Local-packages/math-commands}





% ------------------------------------------------------
% Package "hyperref"
% They advised to put it on the end of preambule
% ------------------------------------------------------
% It allows you to use hyperlinks in the text
\usepackage{hyperref}










% ------------------------------------------------------------------------------------------------------------------
% Title and author of the text
\title{Wstęp do~teorii pola, v0.1}

\author{Na podstawie notatek Andrzeja Herdegena}


% \date{}
% ------------------------------------------------------------------------------------------------------------------










% ####################################################################
\begin{document}
% ####################################################################





% ######################################
\maketitle
% ######################################






% ######################################
\section{Wstępne informacje}
% ######################################


Poniższe notatki powstały na podstawie spotkań na których przerabiane
zostały notatki sporządzone przez prof. Andrzeja Herdegena do jego przedmiotu
„Wstęp do teorii pola”, który pierwszy raz, wedle wiedzy piszącego,
odbył~się w~roku akademickim 2011/2012. Prof. Herdegen zgodził~się
udostępnić swoje notatki zainteresowanym, można je otrzymać choćby od
piszącego te słowa, pisząc pod adres \email. Wszelkie uwagi do tych notatek
i~znalezione w~nich błędy oraz niedoróbki również można pod ten
adres.

Najnowsza wersja tych notatek jest dostępna na GitHubie w~repozytorium
\href{https://github.com/KZiemian/Physics-Learning-materials}
{\texttt{https://github.com/KZiemian/Physics-Learning-materials}}, w~katalogu
\texttt{Wstęp do teorii pola}\footnote{Bezpośredni link: \\
  \href{https://github.com/KZiemian/Physics-Learning-materials/tree/main/Wstęp-do-teorii-pola-Notatki}
{\texttt{https://github.com/KZiemian/Physics-Learning-materials/tree/main/Wstęp-do-teorii-pola-Notatki}}.}, na gałęzi o~nazwie \texttt{Gałąź-robocza}.
Znalezione błędy i~uwagi oczywiście można tworzyć zgłaszać na stronie tego
repozytorium, standardowymi dla GitHuba metodami.

Do sporządzania tych notatek poza notatkami prof. Herdegena korzystano
z~następujących materiałów. Dyskusja czasoprzestrzeni Minkowskiego w~rozdziale \eqref{???} pochodzi z~artykułu prof. Herdegena
\textit{Flat spacetime in a~capsule}
(\parencite{Herdegen-Flat-spacetime-in-a-capsule-Pub-2009}), przy czym
dyskusja przekształceń Lorentza bazuje po części na rozdziale ??? jego
książki \textit{Algebra liniowa i~geometria}
\parencite{Herdegen-Algebra-liniowa-i-geometria-Pub-2010}. Materiał
o~przestrzeniach afinicznych, kwadrykach i~iloczynach tensorowych oraz
zewnętrznych, bazuje na rozdziałach VI, IX i~IV książki Jacka Gancarzewicza
\textit{Algebra liniowa i~jej zastosowania}
\parencite{Gancarzewicz-Algebra-liniowa-i-jej-zastosowania-ETC-Pub-2004}.
Pewne informacje o~problemach sformułowania pełnej teorii pola
elektromagnetycznego i~elektronu zaczerpnięto z~książki Fritza Rohrlicha
\textit{Klasyczna teoria cząstek naładowanych}
\parencite{Rohrlich-Klasyczna-teoria-czastek-naladowanych-Pub-1981}.
























% ####################################################################
% ####################################################################
% Bibliography

\printbibliography





% ############################
% End of the document

\end{document}
